% Folder 149, batch 07 — archivists' pages 121–135.
% Pass: Opus 5.5 (claude-opus-5-5), 2026-10-03 — first pass, unchecked against the pages by a human.
\documentclass[11pt,a4paper]{article}
\input{../preamble/grothendieck}

\folder{149}
\batch{7}
\pages{121}{135}
\foldertitle{[Autour de La "Longue Marche" à travers la théorie de Galois, pages 1 à 67] : notes manuscrites (s.d.).}
\dating{s.d. — le groupe « Autour de La "Longue Marche" à travers la théorie de Galois » (141 à 149) est daté [à partir de 1978]-1983}
\watermark{Édition de démonstration}

\begin{document}

\page{121}
\note{la page commence au milieu d'une phrase : la définition de l'espace $\mathfrak{X}$ des points (« modèle » de $L/K$) et des anneaux locaux $\underline{O}_{\mathfrak{X},x}$ est dans le lot précédent.}

correspondant exactement (via $x \mapsto \underline{O}_{\mathfrak{X},x}$) aux \uncertain{jauges} \uncertain{dans} $L$ qui sont ent.\ de type fini sur $K$, de corps des fractions $L$ — i.e.\ aux $V$ précédents.

\uncertain{Sauf erreur}, \uncertain{ces deux} types de p\supplied{oin}ts de $\mathfrak{X}$ (qui épuisent \struck{tout} si $n = 1$, \uncertain{qui est le seul cas}, \uncertain{avec} $n = 0$, où $\mathfrak{X}$ soit un schéma !) sont les seuls pour lesquels $\underline{O}_{\mathfrak{X},x}$ soit noethérien. Cependant, on a

\marginal{On pourrait \ill{} \uncertain{relation} \ill{} et alors \ill{} \uncertain{relat.} \ill{} $\mathfrak{X}$ \ill{} \uncertain{dessus}…}
\note{note marginale oblique, à gauche de la proposition, reliée au texte par une accolade ; presque illisible.}

\textbf{Proposition.} Pour tout $x \in \mathfrak{X}$, $\underline{O}_{\mathfrak{X},x}$ est un anneau de valuation dans $L$, de corps des fractions $L$, et contenant $K$. L'application $x \mapsto \underline{O}_{\mathfrak{X},x}$ est une bijection de $\mathfrak{X}$ sur l'ens.\ des \uncertain{anneaux} \ill{} \uncertain{anneaux} \uncertain{dans} $L$.

\emph{Dém.} a) Soit $f \in L \setminus \underline{O}_{\mathfrak{X},x}$, il faut prouver $1/f \in \underline{O}_{\mathfrak{X},x}$. Or $f = g/h$, $g, h \in \underline{O}_{X,a}$, où $X$ modèle \add{\uncertain{régulier}} (projectif), $a$ un p\supplied{oin}t, $f$, $g$ \struck{\uncertain{premiers entre eux}} \ill{} — utilisant la résolution, OPS $\exists$ \uncertain{syst.} \struck{\ill{}} \ill{} \struck{\ill{}}, \ill{} $f$, $g$ \uncertain{des} \ill{} de \uncertain{paramètres}, \ill{}

\page{122}
\note{p.~61 de l'auteur.}

un \uncertain{sys.} de paramètres, et on se ramène au cas où $f$, $g$ premiers entre eux, — on peut — faisant encore éclater — \uncertain{supposer} que les diviseurs de $f$, $g$ \struck{\ill{}} \uncertain{se séparent}…

b) Un anneau de val.\ $V$ \struck{qui} pour l'extension $L/K$ domine un p\supplied{oin}t bien déterminé \add{$x$} de chaque $X$, d'où un $x \in \mathfrak{X} = \varprojlim X$, et $V$ domine $\underline{O}_{\mathfrak{X},x}$, donc lui est égal puisque \ill{} \uncertain{dominant} \ill{} \uncertain{est un} \ill{} de valuation.

\textbf{NB} Tout \add{\uncertain{modèle} $U$ de $L/K$ \uncertain{définit un}} ouvert\supplied{s} de $\mathfrak{X}$ \struck{correspondant} (\struck{\ill{}} \uncertain{plonge} dans un modèle propre $X$, c'est clair) et \struck{et} les modèles \uncertain{conven}- \ill{} à \emph{certains} ouverts de $\mathfrak{X}$ (pas stable par réunion finie).

Pour tt anneau local \add{\uncertain{normal}} $\underline{O} \subset L$ de corps des fractions $L$, on peut regarder $\operatorname{Spec} \underline{O} = S$, son p\supplied{oin}t fermé $s$, \uncertain{le normalisé}-

\page{123}
$\underline{O}'$ de $\underline{O}$ dans $\overline{L}$ i.e.\ $S'$ de $S$ dans $\overline{L}$, et pour $x'$ au-dessus de $x$, un groupe de décomposition $N_{x'}$, d'inertie $I_{x'}$, \uncertain{mais} les interprétations \ill{} qu'on a \uncertain{signalé} plus haut dans le cas particulier d'un $\underline{O}$ qui est un anneau jauge. L'idée, c'est que \uncertain{pour} \uncertain{prenant} $\underline{O}$ de valuation (par ex.\ un \uncertain{jauge}-$\underline{O}$) $\supset K$, on trouvera des groupes d'inertie à structure simple — \uncertain{des} $T_{\infty}^{\alpha}$ essentiellement, où $\alpha$ est la « codim » du $x \in \mathfrak{X}$ correspondant —

\marginal{$T_{\infty} = T_{\infty}(\overline{L})$}

et si $x$ est de dim $0$, \uncertain{on a} $\mathcal{E}_{k(x)}^{\overline{k(x)}}$ — \ill{} une extension de $\mathcal{E}_{k(x)}$ — \uncertain{sous}-groupe ouvert de $\mathcal{E}_K$, puisque $k(x)$ est \add{fini} alg.\ sur $K$ — par $T^{\alpha}$, qui \ill{} avec un choix de \uncertain{paramètres} \uncertain{scindées} \uncertain{sous l'effet} du modèle. Ceci donne-

\marginal{on a \ill{} $W$ \uncertain{nécessaires} \ill{} ici}

\page{124}
\note{p.~62 de l'auteur.}

-rait un principe général pour la construction des germes de sections de $\mathcal{E}_L$ sur $\mathcal{E}_K$ — à l'aide des p\supplied{oin}ts de dim $0$ de la « variété de Riemann » $\mathfrak{X}$ de $L$ sur $K$.

Comparant, pour une sous-extension $L_1$ de $K$ de degré de tr.\ $1$, avec l'hom.\ $\mathcal{E}_L \to \mathcal{E}_{L_1}$, on a donc aussi un moyen de construction des \add{germes de} sections de $\mathcal{E}_{L_1}$ sur $\mathcal{E}_K$, et le test intéressant serait de voir si on \add{en} trouve \uncertain{pas} plus que les sections déjà envisagées, associées aux « places » de $L_1$ sur $K$.

\marginal{\uncertain{on pourrait ici} $K = \mathbb{Q}$, \uncertain{au cas fini} \uncertain{typique} !}

Mais il en est bien ainsi, comme on voit en considérant $V_1 = V \cap L_1$, qui est un anneau de valuation de $L_1$,

\page{125}
dont le corps résiduel (contenu dans celui de $V$) sera \struck{alg} fini sur $K$ si celui de $V$ l'est (et $V_1 \neq L_1$) — donc $V_1$ \add{\uncertain{correspond à une}} « place » du corps de fonctions d'une variété $L_1$ sur $K$. Or, par des fonctorialités \uncertain{essentiellement} évidentes, on voit que l'hom $\mathcal{E}_L^{\overline{L}} \to \mathcal{E}_{L_1}^{\overline{L}_1}$ envoie $N_V$ dans $N_{V_1}$, $I_V$ dans $I_{V_1}$ — donc l'image de $N_V$ \add{germe de} normalise $I_{V_1}$, donc (si $\ldots$ une section de $N_V$ sur $\mathcal{E}_K$) le germe de section \uncertain{obtenu} de $\mathcal{E}_{L_1}$ sur $\mathcal{E}_K$ normalise $I_{V_1}$ — et \uncertain{sa} restriction à $\mathcal{E}_K^{\circ}$ centralise $I_{V_1}$ —

\marginal{\ill{} fait voir \uncertain{que} le corps \uncertain{résid.} de $V_1$ \ill{} \uncertain{extension} $L_1$ \ill{} \uncertain{dim} \ill{} \uncertain{fait} \ill{}…}
\note{note marginale oblique en haut à gauche, presque illisible.}

\textbf{Conjecture.} Soient $K$, $L$ des \uncertain{ext.}\ de type fini de $\mathbb{Q}$, $K \subset L$. Alors

a) Toute section de $\mathcal{E}_L$ sur $\mathcal{E}_K$ (ou germe de tel, sa variante $\ldots = \ldots$) normalise un $I_V$ associé à un anneau de valuation $V$ de $L$ contenant $K$, à corps résiduel algébrique \add{sur $K$}, et $V$ est uniquement

\page{126}
\note{p.~63 de l'auteur.}

par cette condition — en fait la section provient d'une section du groupe de décomposition $N_V$ au-dessus de $\mathcal{E}_K$.

b) Soit $U$ un modèle « élémentaire » de $L$ sur $K$, \uncertain{arbitraire}. Alors \uncertain{toute} section de $\mathcal{E}_K^{\overline{L}}$ sur $\mathcal{E}_K^{\overline{K}}$ se relève en une section de $\mathcal{E}_L^{\overline{L}}$ sur $\mathcal{E}_K^{\overline{K}}$.

\marginal{\uncertain{nouveaux pour la question} $N_V = \operatorname{Norm}(\overline{I}_V)$ \uncertain{on voudrait que} $\mathcal{E}_{K_1}^{\overline{K}}$ \ill{} germe de \uncertain{sections} de $\mathcal{E}_L^{\overline{L}}$ \uncertain{associés aux} \ill{} \uncertain{dans} $\overline{L}$ \ill{} de $\mathfrak{X}_{L/K}$ \uncertain{formant un} \ill{} \uncertain{de} \ill{} $I_V$ \uncertain{la} \ill{} \uncertain{sections} \ill{} $H^1(\mathcal{E}_K, \ldots)$ \ill{} \uncertain{associée} \ill{} $\mathcal{E}_K$ \ill{}}
\note{longue note marginale oblique, en deux blocs, le long du bord gauche de la page ; seuls des fragments se lisent. Le second bloc semble porter sur $N_V$, $I_V$, la « formule » \ill{} et $(\mathcal{E}_L^{\overline{L}})^{\Gamma^{\circ}}$.}

Elle peut se \uncertain{décrire} en termes d'un anneau de valuation $V$ de $L/K$, et par $V$ \uncertain{un} centre sur $U$, \uncertain{tel} que se \uncertain{concentre} \uncertain{pour} le cas d'un p\supplied{oin}t $x \in U$ rationnel sur $K$, et de la section définie par $x$. Dans ce cas (posant $\Gamma = \mathcal{E}_K^{\overline{K}}$) on conjecture encore, bien sûr
\[
  (32) \qquad (\mathcal{E}_L^{\overline{L}})^{\Gamma^{\circ}} = (1)
\]
(et on \uncertain{peut} y remplacer $\Gamma^{\circ}$ par $\Gamma^{\circ}_{\ill{}}$ — mais en fait ceci (on \uncertain{voit} ou $\operatorname{Centr}_{\Gamma}(\Gamma^{\circ}) = (1)$) équivaut à la relation analogue \struck{$\pi_{L/K}^{\overline{L}}$}
\[
  (33) \qquad (\pi_{L/K}^{\overline{L}})^{\Gamma^{\circ}} = (1)
\]
qui se déduit d'ailleurs, par dévissage,

\page{127}
de la conjecture analogue déjà faite pour la dim relative $1$.

Pour mieux appréhender cette dernière \add{b)}, partie de la conjecture, il faudrait examiner de plus près ce que deviennent, par composition avec $\mathcal{E}_L^{\overline{L}} \to \mathcal{E}_U^{\overline{L}}$, les sections de $\mathcal{E}_L^{\overline{L}}$ sur $\mathcal{E}_K$ associées à un $V$, et pour commencer, étudier le \add{$\Delta_V$,} sous-groupe de $\pi_{U/K}^{\overline{L}}$ image de $\overline{I}_V$ — qui, en le \uncertain{devine}, \uncertain{donne} une \uncertain{réalisation} algébrique \uncertain{de} \ill{} à un groupe-lacet. Dans quelle mesure des $V$ distincts définissent-elles des $\Delta_V$ distinctes ?

Pour étudier cette question, introduisons un modèle $X$ de $L/K$ contenant $U$, $X$ propre, les comp.\ irréd.\ \add{$D_i$ ($i \in I$)} de \add{$D =$} $X \setminus U$ des diviseurs réguliers, et $X \setminus U$ à croisements normaux. Pour $x \in \operatorname{Supp} D$, considérons le hensélisé \add{(\uncertain{ou strict})} $\widetilde{\underline{O}}_x$ de $\underline{O}_{X,x}$, et l'image inverse $\widetilde{U}_x$ de $U$ dessus. On a

\page{128}
\note{p.~64 de l'auteur.}

$\overline{L}$ au-dessus de $\widetilde{U}_x$, bien sûr, et
\[
  (34) \qquad 1 \to T^{I(x)} \to \mathcal{E}_{\widetilde{U}_x}^{\overline{L}} \longrightarrow \mathcal{E}_{k(x)}^{\overline{k(x)}} \to 1
\]
où \struck{$I(x)$}
\[
  I(x) = \{\, i \in I \mid x \in D_i \,\}, \qquad T = T_{\infty}(\overline{L}^{*}),
\]
$\overline{k(x)}$ étant la clôture alg.\ de $k(x)$ induite par un plongement choisi de $\widetilde{K}_x$ (corps des fractions de $\widetilde{\underline{O}}_x$) dans $\overline{L}$, qui fait donc fonction de clôture sép.\ de $\widetilde{K}_x$.

Ceci dit, pour $V$ anneau de valuation de $L/K$, si son centre $x$ \add{sur $X$} est $\in \operatorname{Supp} D$, on a, \add{(\uncertain{hensélisé})} posant $S = \operatorname{Spec} V$, d'où $\widetilde{S}$ (de p\supplied{oin}t générique $\widetilde{\eta}$), une factorisation de $\widetilde{\eta} \to U$ (compatible $\widetilde{\eta} \to \eta \in \operatorname{Spec}(L)$ $\to U$) en $\widetilde{\eta} \to \widetilde{U}_x$, car $\widetilde{\eta} \to U$ est induit par $\widetilde{S} \to X$ qui se factorise en $\widetilde{S} \to \widetilde{T} \to X$ ($\widetilde{T} = \operatorname{Spec} \widetilde{\underline{O}}_x$). Donc \struck{l'image $N_V$ de} \add{l'hom.\ induit} $\mathcal{E}_{\widetilde{\eta}}^{\overline{L}} \to \mathcal{E}_U^{\overline{L}}$ est \struck{contenue dans celle de} \add{un composé}
\[
  (35) \qquad
  \underset{\substack{\parallel \\ N_V}}{\mathcal{E}_{\widetilde{\eta}}^{\overline{L}}}
  \longrightarrow
  \underset{\substack{\parallel\ \text{dfn} \\ N_x}}{\mathcal{E}_{\widetilde{U}_x}^{\overline{L}}}
  \longrightarrow
  \mathcal{E}_U^{\overline{L}}
\]
(l'image de $\mathcal{E}_{\widetilde{U}_x}^{\overline{L}} \to \mathcal{E}_U^{\overline{L}}$ n'étant d'ailleurs autre que le groupe de décomposition de $x$ dans $\mathcal{E}_U^{\overline{L}}$ relatif à $\overline{L}$…). L'étude de l'hom composé
\note{sous le second terme de (35), un symbole biffé, illisible, avant « $N_x$ ».}

\page{129}
se ramène à celle des hom composants. À ce sujet, sous réserve de vérifications, je présume

1°) que l'hom $\mathcal{E}_{\widetilde{U}_x}^{\overline{L}} \to \mathcal{E}_U^{\overline{L}}$ est \emph{injectif} — ou ce qui revient au \uncertain{même}, que \uncertain{son} partie géométrique
\[
  \underset{\substack{\parallel \\ I_x \simeq T^{(I(x))}}}{\mathcal{E}_{\widetilde{U}_x/K}^{\overline{L}}} \longrightarrow \pi_{U/K}^{\overline{L}}
\]
est \emph{injective}.

\marginal{\uncertain{les deux lettres} $I$ \uncertain{ne sont pas les mêmes}}
\marginal{C'est \uncertain{sûrement} « \uncertain{en général} » \emph{faux} \add{pour $\operatorname{card} I(x) \neq 1$} comme on voit en \uncertain{remplaçant} $X$ par un autre $X'$ dominant $X$, \uncertain{en éclatant} des \ill{} \uncertain{faisant} un p\supplied{oin}t de $D$… \textbf{Mais} \uncertain{est-il vrai} \uncertain{qu'on} \ill{} $X \supset U$ \uncertain{tel que} \ill{} \uncertain{localement} \emph{injectif}…}
\note{note marginale oblique dans la marge droite, d'une encre plus foncée, vraisemblablement ajoutée après coup ; elle porte sur l'énoncé 1°).}

2°) que \struck{l'hom} lorsque le corps \uncertain{résiduel} \ill{} $V$ est alg$/K$ (i.e.\ $V$ correspond \uncertain{à un point} de $\mathfrak{X}_{L/K}$ de dim zéro), alors l'hom
\[
  \mathcal{E}_{\widetilde{\eta}}^{\overline{L}} = N_V \to \mathcal{E}_{\widetilde{U}_x}^{\overline{L}}
\]
\struck{est égal} \add{a une image} \struck{surjectif} \add{ouverte} — ce qui revient à dire (puisque $k(V)$ \struck{alg} \add{fini} sur $k(x)$ donc $\mathcal{E}_{k(V)} = N_V / I_V \hookrightarrow \mathcal{E}_{k(x)}$ \add{(sous-groupe ouv.)} $\simeq N_x / I_x$), que la partie géométrique
\[
  I_V \longrightarrow I_x = T^{I(x)}
\]
\struck{soit \uncertain{à image} \ill{} \uncertain{ouverte}} — \struck{\uncertain{ou plus} \ill{}} — \add{\uncertain{soit}} surjective. \struck{\ill{} \uncertain{sous-groupe ouvert}.}
\note{la fin de 2°) a été biffée et récrite plusieurs fois ; seul « surjective » est sûr.}

À noter que la question 2° est purement

\page{130}
\note{p.~65 de l'auteur.}

locale autour de $x$ — elle doit être essentiellement « triviale », que ma présomption s'avère vraie ou fausse — par contre 1°) est une question de nature globale sur $U$, et sans doute \struck{d'autre} \uncertain{aucunement} façon triviale.

\struck{Avant de \uncertain{l'oser}} La validité de ces énoncés \struck{signifierait} \add{impliquerait} donc, pour les sections de $\mathcal{E}_U^{\overline{L}}$ sur $\mathcal{E}_K^{\overline{L}}$ associées à un \struck{\ill{}} \add{anneau} de valuation de $L/K$, \add{de corps résiduel $K$} que l'image de $\mathcal{E}_K^{\overline{L}}$ doit normaliser un \struck{$I_x$} sous-groupe
\[
  \Delta_x \simeq T^{I(x)} \text{ de } \pi_{L/K}^{\overline{L}},
\]
qui est non trivial si la valuation n'\uncertain{a} pas de centre sur $U$, i.e.\ si la section n'est pas associée à un p\supplied{oin}t $K$-rationnel de $U$, ce qui est justement dans le sens des conjectures (qui pouvaient d'abord sembler hérétiques !) du §~2.
\note{le §~2 auquel renvoie le texte est antérieur à ce lot.}

Avant de me tourner vers l'étude des questions 1°) et 2°) (et \add{aussi} des questions de normalisation et de centralisation soulevées

\page{131}
précédemment à propos de $N_V$, $I_V$ …), je vais regarder d'un peu plus près l'allure des sous-groupes $I_x \subset \pi_{L/K}^{\overline{L}}$, \struck{associés} \uncertain{ou plutôt}, des classes de \add{$\pi$-}conjugaison de tels sous-groupes) associés aux $x \in X \setminus U$, \add{pour $x$ variable,} (et une connaissance de caractère intrinsèque de ces groupes, indépendamment du choix du modèle $X$ (mais ceci est évident, si on admet ce qui vient d'\uncertain{être} \uncertain{présumé} — \uncertain{que ce sont justement}, à \uncertain{conjugaison} — et \uncertain{à} \uncertain{sous-groupes} ouverts près \uncertain{tout} au plus, les images des $I_V$, pour $V \in \mathfrak{X}_{L/K}$ \struck{\ill{}} \uncertain{n'ayant} pas de centre dans \struck{$X \setminus U$} $U$ …)

\struck{Considérons d'abord les}

Soit $J$ une partie de \struck{l'}l'ens.\ fini $I$, et soit
\[
  D_J^{*} = \{\, x \in X \mid I(x) = J \,\} = \overbrace{\bigcap_{j \in J} D_j}^{D_J} - \bigcup_{i \in I \setminus J} D_J \cap D_i ,
\]
\[
  D_{\{i\}} = D_i \setminus \bigcup_{j \neq i} (D_i \cap D_j).
\]
Si $\alpha = \operatorname{card} J$, on sait que $D_J \subset X$ est \add{régulier,} de codim $\alpha$ en chacun de ses points (mais bien sûr, il peut être $\emptyset$). Si on appelle « simplexe » une partie $J \subset I$

\marginal{il \uncertain{vient} \uncertain{un nouveau} \uncertain{décrire} $D_J$ \uncertain{comme} \uncertain{l'ensemble} $\bigcap_{i \in J} D_i$ (36) \uncertain{et} \uncertain{décrire} \uncertain{comme} $\{x \in X \mid I(x) \supset J\}$}
\note{note marginale oblique, en bas à gauche ; lecture très incertaine. Dans la seconde formule, $D_{\{i\}}$ est écrit sans astérisque ; il s'agit vraisemblablement de $D^{*}_{\{i\}}$.}

\page{132}
\note{p.~66 de l'auteur.}

telle que $D_J \neq \emptyset$ \add{(i.e.\ $D_J = \bigcap_{j \in J} D_j$ soit $\neq \emptyset$)}, on trouve, sur $I$ une structure de \struck{schéma} \add{\uncertain{ensemble}} simplicial — quelle est sa signification topologique ? Les composantes connexes correspondent \supplied{aux} comp.\ connexes de $D = X \setminus U$ — on connaît : celles du \uncertain{bout} de $\widehat{X} \setminus D$, où $\widehat{X}$ est le \add{\uncertain{formel}} complété de $X$ le long de $D$ — et $\widehat{X} \setminus D$ \ill{} une \uncertain{notion} \uncertain{invariante} \uncertain{p.r.}\ à la « complétification » choisie $X$ de $U$. S'il était vrai que les $D_J$ \add{non vides} sont $\infty$-connexes (ce qui bien sûr n'est pas le cas, si $J$ \uncertain{pour} $D_J$ \uncertain{réduit} à un point !) on pourrait penser que l'\struck{schéma} \add{\uncertain{ensemble}} simplicial exprime la topologie du type d'homotopie de $U_{\infty}$.

Considérons maintenant le « localisé » \struck{\ill{}} (en un sens \uncertain{à préciser} convenable) de $X$ le long de $D_J$, soit \struck{$X$} $X_J$, et soit $U_J$ \add{$= X_J \setminus D_J$} la trace de $U$ sur $X_J$. Considérons $\mathcal{E}_{U_J}$ (de façon précise, $\mathcal{E}_U^{\overline{L}}$, si on localise par \uncertain{p.v.}\ algébrique), on trouve une suite exacte
\[
  (37) \qquad T^{J} \to \mathcal{E}_{U_J}^{\overline{L}} \to \mathcal{E}_{D_J} \to 1
\]
— de façon plus précise, choisissons un rev.\ universel $\widetilde{D}_J$ de $D_J$, on trouve un rev.\ universel (par image inverse \add{d'ailleurs}) $\widetilde{X}_J$ de $X_J$, d'où \uncertain{aussi} $\widetilde{U}_J$ et $\widetilde{U}_J \to U$ ; \struck{et une} on désigne $\widetilde{L}_J$ le corps des fonctions \add{algébr.}\ de $\widetilde{X}_J - \widetilde{U}_J$, c'est une ext.\ \add{de $L$}, et

\marginal{il \uncertain{est} \uncertain{plus} \uncertain{convenable} \uncertain{de prendre} les « \uncertain{tubulaires} » \ill{} \uncertain{ou} « \uncertain{formels} »…}
\marginal{\ill{} \uncertain{dessus} $\widetilde{D}_J$ \ill{} $U_J$ \ill{} \uncertain{revêtement universel} \ill{} $D_J$ \ill{}}
\note{deux notes marginales obliques à gauche, la seconde très serrée en bas de page ; seuls des fragments se lisent. Au-dessus de $\mathcal{E}_{D_J}$ dans (37), un exposant biffé ($\widetilde{D}_J$ ?).}

\page{133}
et on choisit un \add{$L$-}plongement de $\widetilde{L}_J$ dans $\overline{L}$. C'est par cela que (37) est définie — c'est la \add{\uncertain{début} de la suite d'homotopie du} fibré, \uncertain{morale}, d'un produit de disques \struck{\ill{}} sur $D_J$ (\uncertain{savoir} $U_J$ \uncertain{homotope} $=$ ) — j'ai eu un peu peur d'y mettre \struck{\uncertain{les}} \uncertain{zéro} à \uncertain{gauche}, à cause du \struck{\ill{}}
\[
  \pi_2(\overline{D}_J) \to \pi_1(\text{fibre}) \qquad (?)\,!
\]
\note{sous « $\pi_1(\text{fibre})$ » est écrit « $= T^J$ ».}

Mais considérons l'hom
\[
  \mathcal{E}_{U_J}^{\widetilde{U}_J} \longrightarrow \mathcal{E}_U^{\overline{L}}
\]
induit par « l'injection » $U_J \to U$, \struck{\uncertain{elle}} heuristiquement on \struck{devrait avoir} \add{\uncertain{devrait}} un diagramme \struck{induit par} commutatif

\begin{tikzcd}
U_J \arrow[r] \arrow[d, "i"'] & U \arrow[d] \\
X_J \arrow[r] & \operatorname{Spec} K = \xi \\
D_J \arrow[u, "\text{équiv.\ d'homotopie}"']
\end{tikzcd}

d'où un hom.\ de suites exactes d'homotopie — \add{de cette relation} \struck{relative} à la première \struck{\uncertain{ligne}} \add{\uncertain{cohomologique}}, dans cette relation : $T^J$ \uncertain{scindée}…

\begin{tikzcd}[column sep=small, row sep=small, nodes={font=\scriptsize}]
 & \pi_2(D_J) & T^{J} & \mathcal{E}_{U_J} & \mathcal{E}_{D_J} & \\
\cdots \arrow[r] & \pi_2(X_J) \arrow[u, no head, "\parallel"'] \arrow[r] \arrow[d] & \pi_1(\overline{\mathbb{G}}_m^{\,J}) \arrow[u, no head, "\wr"'] \arrow[r] \arrow[d] & \pi_1(U_J) \arrow[u, no head, "\wr"'] \arrow[r] \arrow[d] & \pi_1(X_J) \arrow[u, no head, "\wr"'] \arrow[r] \arrow[d] & 1 \\
\cdots \arrow[r] & \pi_2(\xi) \arrow[r] \arrow[d, no head, "\parallel"'] & \pi_1(U_{\overline{\xi}}) \arrow[r] \arrow[d, no head, "\wr"'] & \pi_1(U) \arrow[r] \arrow[d, no head, "\wr"'] & \pi_1(\xi) \arrow[r] \arrow[d, no head, "\wr"'] & 1 \\
 & 0 & \pi_{U/K} & \mathcal{E}_U & \mathcal{E}_K &
\end{tikzcd}

qui montre que \add{le composé} $\pi_2(D_J) \to T^J \to \pi_{U/K}$ est nul — donc si on admet que $T^J \to \pi_{U/K}$ est injectif, on conclura que $\pi_2(D_J) \to T^J$ est nul !

\page{134}
\note{p.~67 de l'auteur.}

On va donc admettre que (37) a, \uncertain{définit} bien une structure d'extensions, et on a \uncertain{aussi} un hom de suites exactes

\note{à gauche du diagramme, deux numéros : un « (38') » biffé ou surchargé sur la première ligne, et « (38) » sur la seconde.}

\begin{tikzcd}[column sep=small]
1 \arrow[r] & T^{J} \arrow[r] \arrow[d] & \mathcal{E}_{U_{J,\alpha}}^{\overline{L}} \arrow[r] \arrow[d] & \mathcal{E}_{D_{J,\alpha}}^{\widetilde{D}_{J,\alpha}} \arrow[r] \arrow[d] & 1 \\
1 \arrow[r] & \pi_{U/K}^{\overline{L}} \arrow[r] & \mathcal{E}_U^{\overline{L}} \arrow[r] & \mathcal{E}_K^{\overline{K}} \arrow[r] & 1
\end{tikzcd}

— où, pour \uncertain{simplifier} la \uncertain{diction}, on \uncertain{fait} \uncertain{abstraction} \uncertain{du} \uncertain{fait} que (37) \ill{} de sens telle quelle que si \add{i.e.\ $X_J$ connexe} $D_J$ est connexe, que \uncertain{sinon} il faut regarder les composantes connexes $X_{J,\alpha}$ de $X_J$, \uncertain{induisant} les composantes conn.\ $D_{J,\alpha}$ de $D_J$, et $U_{J,\alpha}$ de $U_J$, et écrire la suite exacte pour chaque composante. Bien sûr, il y a des choses plus intrinsèques : ici \uncertain{avec} les \uncertain{types} \struck{et les \uncertain{groupoïdes}} \add{multiplicatifs}, \uncertain{fondamentaux}… et leurs groupoïdes \struck{classes de $\mathcal{E}_U^{\overline{L}}$-conjugaison}

On voit ainsi que les hom
\[
  (39) \qquad T^{J} \longrightarrow \pi_{U/K}^{\overline{L}}
\]
induits par les différents $x \in D_J$ — chacun n'étant en fait, quand on reprend la définition précise, que défini modulo opération d'automorphismes intérieurs de $\mathcal{E}_U^{\overline{L}}$, induits par

\page{135}
\struck{\ill{}} éléments dont l'image dans $\mathcal{E}_K^{\overline{K}}$ est dans l'image de $\mathcal{E}_{k(x)}^{\overline{k(x)}} \to \mathcal{E}_K^{\overline{K}}$ — \uncertain{ont} que pour $x$ variable dans une composante connexe $D_{J,\alpha}$ de $D_J$, sont tous conjugués entre eux dans \struck{\ill{}} $\mathcal{E}_{U/K}$ — et on peut faire les choix relatifs aux différents $x$ d'une façon cohérente, qui mainten\supplied{ant} l'incertitude extérieure de \uncertain{celle}, qui exprime l'indétermination dans $\pi_{U/K}^{\overline{L}}$ \uncertain{ici} \uncertain{seulement} aux produits des éléments de l'image de
\[
  (39) \qquad \mathcal{E}_{D_{J,\alpha}}^{\widetilde{D}_{J,\alpha}} \longrightarrow \mathcal{E}_K^{\overline{K}} .
\]
$J = \{i\}$ est particulier, \struck{\uncertain{lorsque}} \add{dans} ce cas d'ailleurs $D_J = D_i$ \struck{\uncertain{composante}}, \add{et} \uncertain{connexe}, et on trouve des hom
\[
  (40) \qquad T \xrightarrow{\ \kappa_i\ } \pi_{U/K}^{\overline{L}} ,
\]
déterminés a priori mod.\ automorphismes intérieurs \uncertain{\ill{}} — dans un \uncertain{fait} $\kappa_i$ est un hom \emph{extérieur} bien déterminé, qui est \struck{\ill{}} compatible avec les opérations \struck{\ill{} par \uncertain{opération}} \uncertain{ext.}\ de $\mathcal{E}_K^{\overline{K}}$ sur $T$, et sur $\pi_{U/K}^{\overline{L}}$.

\struck{Si $J' \supset J$ (i.e.\ $D_{J'} \subset D_J$), \uncertain{choisissons} une composante conn.\ $D_{J',\alpha'}$ \ill{} de $D_{J'}$, elle est contenue dans une \uncertain{seule} comp.\ conn.}
\note{ce dernier paragraphe (avec une variante biffée « (i.e.\ $D_{J'} \neq \emptyset$) ») est barré de grandes croix ; la page — et le lot — s'arrêtent là, la suite est dans le lot suivant. Le numéro (39) est porté ici une seconde fois, déjà employé p.~134.}

\marginal{\uncertain{de façon plus précise}, le choix \uncertain{d'un} \uncertain{revêtement universel} \uncertain{en impliqu}\supplied{ant} $\widetilde{U}_x$ \ill{} \uncertain{dans} \ill{} $\overline{K}$ \ill{} \uncertain{définition} \ill{} \uncertain{ainsi} $\widetilde{U}_x$ \ill{} \uncertain{universel} \ill{} \uncertain{faites}, $\mathcal{E}_K^{\overline{K}}$ \ill{} $\pi_{U/K}^{\overline{L}}$ \ill{} \uncertain{connexes} \ill{}}
\note{longue note marginale oblique à gauche, sur toute la hauteur de la page, d'une écriture très serrée ; seuls quelques mots se lisent. Elle paraît préciser les choix de revêtements universels qui rendent les $\kappa_i$ bien définis.}

\end{document}
